% paper/sections/01_introduction.tex
\section{Introduction}
\label{sec:intro}

Large language models (LLMs) have become essential infrastructure for a wide range of applications, but their deployment introduces a fundamental cost-quality tension: frontier models like GPT-4o deliver superior performance but cost $16\times$ more per token than smaller models like GPT-4o-mini. For applications processing millions of queries, routing each query to the appropriate model is critical.

Existing routing approaches fall into two categories. \textbf{Static supervised routers} (e.g., RouteLLM \citep{ong2024routellm}, Not Diamond \citep{notdiamond2024}, FrugalGPT \citep{chen2023frugalgpt}) learn a classifier from large human-preference datasets (55K+ labels for RouteLLM-SW) and apply it at inference time without adaptation. \textbf{Heuristic cascades} (e.g., FrugalGPT, AutoMix \citep{aggarwal2024automix}, Dekoninck et al.\ \citep{dekoninck2025unified}) call a weak model first and escalate based on confidence, but require threshold tuning and cannot handle simultaneous multi-model selection.

Static supervised routers share three structural limitations: (1) they require substantial pretraining data before routing any query, (2) they are architecturally binary or require $O(K^2)$ separate classifiers for $K$ models, and (3) they do not adapt to changing query distributions or model behavior in production. Cascades address (1) but share (2) and (3).

\textbf{Our approach.} We frame LLM routing as a contextual bandit problem and introduce \textbf{RouteSmith}, a multi-strategy routing system. This paper focuses on its online bandit component, which uses two adapted contextual bandit algorithms over a 27-dimensional query feature space: \textbf{LinUCB-27d}, which achieves strong APGR by combining UCB exploration with per-query linguistic features, and \textbf{LinTS-27d}, which applies Linear Thompson Sampling \citep{abeille2017linear} to LLM routing, eliminating the exploration hyperparameter via posterior sampling. The full RouteSmith product (see \S\ref{sec:product}) additionally includes a random forest quality predictor and cold-start $\rightarrow$ warm-up $\rightarrow$ learned phase transitions, described in \citet{routesmith2025system}. Both bandit algorithms require zero pretraining data, learn online from feedback, and extend naturally to $K > 2$ arms\mbox{---}a structural advance over all existing LLM routers.

\textbf{Contributions:}
\begin{enumerate}
    \item \textbf{Contextual bandit routing for LLMs:} Application of LinUCB and LinTS (adapted from \citet{li2010contextual} and \citet{abeille2017linear}) to LLM routing over a 27-dim linguistic feature space. LinUCB-27d achieves APGR=1.126 on MMLU; LinTS-27d (parameter-free) achieves APGR=0.593 with 46\% cost reduction vs. Static-Strong.
    \item \textbf{Binary routing evaluation (Experiment 1):} Comparison on MMLU (600 questions) and GSM8K (300 questions) against Static-Strong, Static-Weak, Random, RouteLLM-SW (three thresholds), and TS-Cat. We note that the RouteLLM-SW comparison uses fallback embeddings and should be interpreted as a domain-transfer stress test, not a head-to-head defeat.
    \item \textbf{Multi-model routing (Experiment 2):} 5-arm online routing across GPT-4o, Claude-Sonnet-4.5, Qwen-Plus, MiniMax-M1, and DeepSeek-V3\mbox{---}structurally impossible for binary routers without $O(K^2)$ re-training.
    \item \textbf{Ablations:} Feature dimensionality, warm-start labels, $\beta$ sensitivity, and $N$-arm convergence analysis.
\end{enumerate}

All results are from real API calls recorded in reproducible JSON result files. No simulated data. We include all hyperparameters in the stored results; formal statistical significance testing is left to future work given the limited number of independent seeds (5).
